\documentclass[landscape]{sciposter}

\usepackage{a0size}
\usepackage{amsmath}
\usepackage{amssymb}
\usepackage{multicol}
\usepackage[english]{babel}
\usepackage{setspace}
\usepackage[pdftex,final]{graphicx} 
\usepackage{epsfig}
\usepackage{dsfont}
\usepackage{bm}
\usepackage{color}
\usepackage{url}


\newcommand{\refeqn}[1]{ (\!\!~\ref{eq:#1}) } % gives references to
\newcommand{\refthm}[1]{ (\!\!~\ref{#1}) }    % equations or theorems





\definecolor{BoxCol}{RGB}{0,120,0 }


\definecolor{SectionCol}{rgb}{1,1,1}
% uncomment for dark blue \section text
\definecolor{TextColG}{rgb}{0,0.5,0}
\definecolor{TextColR}{rgb}{0.5,0,0}
\newcommand{\Rcolor}{\textcolor{TextColR}}
\newcommand{\Gcolor}{\textcolor{TextColG}}

\definecolor{OSUGreen}{RGB}{0,120,0}
\newcommand{\Green}{\textcolor{OSUGreen}}

\definecolor{Gold}{RGB}{255,215,0}
\newcommand{\Gold}{\textcolor{Gold}}

\definecolor{RoyalPurple}{RGB}{120,81,169}
\newcommand{\Purple}{\textcolor{RoyalPurple}}

\definecolor{Black}{RGB}{0,0,0}
\newcommand{\Black}{\textcolor{Black}}


%%%%%%%%%%%%%%%%%%%%%%%%%%%%%%%%%%%%%%%%%%%%%%%%%%%%%%%%%%%%%%%%%%
%My new commands







%%%%%%%%%%%%%%%%%%%%%%%%%%%%%%%%%%%%%%%%%%%%%%%%%%%%%%%%%%%%%%%%%%

%-------------------------------------------------------------------------------------------------------
\title{\Green{The BEST Project Ever!!}}




\author{ \hspace{0 in} Good Student }

\institute{\Black{ \hspace{0 in} State University of New York at Oswego\\ \hspace{0 in} Department of Mathematics\hspace{3.1 in}
\\Oswego, NY, USA\\}}

\email{gstudent3@oswego.edu}

% The following commands can be used to alter the default logo settings
\leftlogo[2]{oswegologo}  % defines logo to left of title (with scale factor)
\rightlogo[1.0]{oswegologo}  % same but on right

%-------------------------------------------------------------------------------------------------------
\begin{document}

%define conference poster is presented at (appears as footer)
\conference{SUNY Oswego Summer Scholar 2014}

\maketitle

% 
\begin{abstract}
\begin{center}
\parbox{35 in}{\Large The abstract should be a brief overview of the project.  Keep specialized notation to a minimum; try to explain to a general audience what the project is about in a way to draw a viewer's attention.  Some say that the abstract of a poster should have a larger font size, to catch a viewer's eye as he/she passes along but if you don't want to have a larger font size then delete the sizing command at the start of this textbox. You can also make your abstract horizontally skinnier by decreasing the number in the parbox qualifier.}
\end{center}
\end{abstract}

%%% Begin of Multicols-Enviroment -- 
\begin{multicols}{4}

%-------------------------------------------------------------------------------------------------------
\section*{INTRODUCTION}
% To leave a section number out, put a star in.

Here's where you give the viewer a reason for wanting to read your poster  \cite{CD}.  How is your project important?  Does it have the potential to make a big impact? What's interesting about your project? 
\\
\\
This is also where you might introduce special notation or specialized vocabulary.  If you can recommend further background reading \cite{Huppert} then you could do that here.  
\\
\\
This section text is currently peppered with citations so that there will be references to put in the References section, and also to give you an idea of how to ask the compiler to provide citations.


%-------------------------------------------------------------------------------------------------------
\section{NOTATION}

Some good rules for mathematical notation:
\\
\\
\begin{itemize}
\item All variables and symbols should be in math mode, surrounded by \$'s so that the {\it symbols} can be distinguished from ordinary {\it text}.
\item Long equations should be displayed either in display math mode or in an equation environment.
\item Overfull and underfull hboxes are annoying for the compiler.  You can try to fix them but sometimes \dots just leave them alone.
\item Notation is supposed to make it easier to understand the math.  Extra notation makes the math more difficult to follow.  It's a fine line between enough notation and too much notation.
\end{itemize}

\

\begin{center}
\includegraphics[scale = 2.6]{oswegologo}
\end{center}

%%-------------------------------------------------------------------------------------------------------

\section{FIGURES}
Blanks are bad for business. Fill your space with figures, if you have nothing else.


\begin{center}
\includegraphics[scale=3.]{oswegologo}
\end{center}

Figures, besides silly logos, should probably have a caption and be numbered.  You can use the figure environment to get numbers and captions.

\begin{figure}
\centering
\caption{Let's put a caption here!}
\includegraphics[scale=2.7]{oswegologo}
\end{figure}

\begin{figure}
\centering
\includegraphics[scale=2.7]{oswegologo}
\caption{No! Let's put the caption here!}
\end{figure}

Don't forget that you may need to scale your images so that they fit in the columns.  An oversized image will wreak havoc on the appearance of your poster.

%%-------------------------------------------------------------------------------------------------------

\section{The Multicols Environment}

The \{4\} after the multicols command tells the compiler to divide the width of the poster into four equal columns, and adjust the length of the poster as necessary.  If you have less content or wide figures then you might be better off with fewer columns -- three is fine, but 2 starts to make things look sparse.
\\
\\
The thing is, you need to have enough material to fill up the poster and prevent a giant white spot at the bottom.  That's just awful-looking!


%-------------------------------------------------------------------------------------------------------
\section{MORE SECTIONS!}

The more we know about our degradation function H and our noise function $N$, the better our estimate of $f(x,y)$ will be. In the spatial domain, we represent our degraded function as: 
\\
\begin{equation}
g(x,y)=h(x,y) \ast f(x,y) + \eta(x,y)
\label{eq:conv}
\end{equation}
\\
where $h(x,y)$ is the spatial representation of the degradation function. 

%-------------------------------------------------------------------------------------------------------
\section{References with BibTeX}

Open the references.bib file and check out the sample citation entries.  If you require different kinds of entries, search ``Bibtex entries'' on the internet and a whole array of options will be available.  Each entry has a name that you'll use in the tex file when you're typing to call up that citation.
\\
\\
Once your references are each formatted as a BibTeX entry, you can cite them using the cite\{\} command.  in the tex tile.  There are several examples here in this document.  If you don't actually have a citation in the text for something, you can use the command ``nocite'' in front of the bibliography style command in your file.  
\\
\\
You need to cite data sources and any specialized software that you've used (GAP, Sage, R).  Usually specialized software programs have information on their webpages about how to properly cite the software.  
\\
\\
Now you need to compile your tex file.  First, make sure references.bib and your tex file are in the same folder.  Then latex compile your tex source file.  If you're using TeXShop or TeXWorks, then the next step is to Bibtex compile your tex source file.  Then latex compile your tex source file TWICE more.  It's crazy!  If you're using other software, sometimes you just need to latex compile your tex source file 3-4 times.  When in doubt, the internet can help you fgure out the procedure.

\begin{figure}
\centering
\includegraphics[scale = 2.0]{oswegologo}
\end{figure}



%---------------------------------------------------------------------------------------------------------------
\section{ACKNOWLEDGEMENTS}

You MUST acknowledge your funding sources, it's polite to acknowledge anyone who has significantly helped you in your project.

%-------------------------------------------------------------------------------------------------------

\nocite{WinNT}
\nocite{GAP4}
\bibliographystyle{plain}
\bibliography{references}

\end{multicols}

% This makes a green bar along the bottom of the poster content.  If you want to add a message, feel free.
\section*{ }

\end{document}
